\section{The Air Freight Substitution Market}
\label{sec:air}

The headline economic case for the 48-hour corridor is simple to state and requires care
to estimate rigorously. The United States domestic air cargo market is approximately
\$65 billion per year \citep{BTS_airfreight2023}. A substantial fraction of that market
consists of goods that fly not because air is the optimal mode but because 87-hour truck
is incompatible with the product's time requirements. If reliable 48-hour truck replaces
air for those goods, the freight cost falls by 90\% or more. The question is how large
that fraction is.

\subsection{The Structure of Domestic Air Cargo}

Bureau of Transportation Statistics data for 2023 record 8.7 billion ton-miles of domestic
air cargo, valued at approximately \$65 billion in freight charges \citep{BTS_airfreight2023}.
The five largest domestic air freight routes by ton-miles are all coast-to-coast or
long-haul connections: Los Angeles -- New York, San Francisco -- New York, Los Angeles
-- Chicago, Seattle -- New York, and Los Angeles -- Miami. Together, these five
city-pairs and their near-parallels account for roughly 23\% of domestic air cargo
ton-miles.

The commodity mix on coast-to-coast air cargo lanes reflects the product categories where
transit time is genuinely binding. FAF5 commodity flow data \citep{FAF5_2023} indicate the
following approximate breakdown for NY--LA and parallel coast-to-coast routes:

\begin{table}[h]
\centering
\caption{Commodity composition of coast-to-coast air cargo (\$15B estimated total)}
\label{tab:air_commodities}
\begin{tabular}{lrrrr}
\toprule
Commodity Class & Share & Est.\ Value & Air \$/lb & Truck \$/lb \\
\midrule
Pharmaceuticals \& biotech        & 30\% & \$4.5B & 8.00--15.00 & 0.80--1.50 \\
Electronics \& semiconductors     & 25\% & \$3.8B & 4.00--8.00  & 0.25--0.50 \\
Fashion \& apparel                & 15\% & \$2.3B & 3.00--6.00  & 0.20--0.40 \\
Perishables (produce, seafood)    & 15\% & \$2.3B & 3.00--5.00  & 0.35--0.60 \\
Industrial parts \& equipment     & 15\% & \$2.3B & 2.00--4.00  & 0.20--0.35 \\
\midrule
\textbf{Total}                    & \textbf{100\%} & \textbf{\$15.1B} & & \\
\bottomrule
\multicolumn{5}{l}{\footnotesize Sources: \citet{BTS_airfreight2023}, \citet{FAF5_2023}. Shares are approximate.} \\
\end{tabular}
\end{table}

The \$15 billion estimate for coast-to-coast air cargo is derived from the BTS corridor
share analysis: the five largest coast-to-coast city pairs account for approximately
23\% of domestic air freight ton-miles, but a higher share of freight value (because
coast-to-coast shipments tend toward high-value goods). We estimate 22--25\% of the
\$65 billion market, or approximately \$14--16 billion, on coast-to-coast corridors
where I2.0 relay would offer a viable alternative.

\subsection{What Stays on the Plane}

Not all air cargo can or should shift to truck, even at 48-hour reliability.
Three categories are excluded from the addressable market:

\textit{Sub-24-hour urgency freight.} Human organs for transplant, blood products,
emergency response materials, and certain time-critical industrial components
(a parts-out airline waiting on an AOG shipment) require transit times under 12 hours.
This category is genuinely time-sensitive and will remain on aircraft regardless of
truck reliability. We estimate this at 10--15\% of coast-to-coast air cargo.

\textit{Ultra-cold-chain pharmaceuticals.} mRNA vaccines and other products requiring
storage at $-70$\,\textdegree C or colder cannot currently be transported in refrigerated
truck containers for 48-hour transits. The cold chain equipment for these products is
designed for short-duration, air-based logistics. This category (perhaps 5\% of air pharma
volume) remains air-dependent regardless of transit time.

\textit{Live animals and specialized cargo.} Live animals, certain hazardous materials,
and cargo requiring in-flight monitoring remain on aircraft. Small category; included
for completeness.

Excluding these categories, approximately 75--80\% of coast-to-coast air cargo is
\textit{technically} substitutable with 48-hour refrigerated truck. Whether shippers
actually substitute depends on service reliability, carrier availability, and shipper
inertia. We apply a conservative 60\% adoption rate, producing an addressable market
of approximately \$9 billion.

\subsection{The Cost Calculation}

The freight cost saving from air-to-truck substitution depends on the price differential
for each commodity class. Air freight averages approximately \$4 per pound on domestic
routes, with higher rates for temperature-sensitive and high-value goods
\citep{BTS_airfreight2023}. Refrigerated truck freight on a 2,790-mile haul costs
approximately \$0.30--0.40 per pound, including the temperature-control premium
\citep{FHWA_freight2023}.

For the \$9 billion addressable market:

\begin{equation}
\text{Freight cost savings} = \$9\text{B} \times \left(1 - \frac{\$0.35}{\$4.00}\right) = \$9\text{B} \times 0.91 = \$8.2\text{B/year}
\end{equation}

\noindent
This is a logistics cost transfer, not a GDP addition. Shippers save \$8.2 billion per year;
airlines lose an equivalent revenue stream from lower-value cargo. The net effect on
GDP depends on how the cost savings propagate through supply chains --- lower final goods
prices, expanded market access for producers, and reduced inventory carrying costs all
represent genuine economic gains beyond the raw logistics cost reduction.

\begin{table}[h]
\centering
\caption{Air-to-truck substitution economics by commodity class}
\label{tab:substitution}
\begin{tabular}{lrrrr}
\toprule
Category & Volume (\$B) & Addressable & Air cost & Savings \\
\midrule
Pharmaceuticals \& biotech     & \$4.5B & 70\% = \$3.2B & \$10/lb  & \$2.9B \\
Electronics \& semiconductors  & \$3.8B & 65\% = \$2.5B & \$5/lb   & \$2.3B \\
Fashion \& apparel             & \$2.3B & 90\% = \$2.1B & \$4/lb   & \$1.9B \\
Perishables                    & \$2.3B & 75\% = \$1.7B & \$4/lb   & \$1.5B \\
Industrial parts               & \$2.3B & 70\% = \$1.6B & \$3/lb   & \$1.4B \\
\midrule
\textbf{Total (60\% adoption)} & \textbf{\$15.1B} & \textbf{\$9.1B} & & \textbf{\$8.2B} \\
\bottomrule
\multicolumn{5}{l}{\footnotesize Truck cost assumed at \$0.35/lb (refrigerated). Sources: \citet{BTS_airfreight2023}, \citet{FHWA_freight2023}.} \\
\end{tabular}
\end{table}

\subsection{The Parallel Corridor Multiplier}

The \$8.2 billion annual savings estimate covers the New York--Los Angeles corridor and
its direct parallels (NY--San Francisco, NY--Seattle, NY--Portland). It excludes the
Chicago--Los Angeles corridor (I-80 west from Chicago), the Boston--Los Angeles corridor,
and the entire Southeast--West Coast market. These parallel corridors would benefit from
I2.0 relay networks on the same infrastructure, at minimal additional cost, because the
relay stations are shared.

Conservative estimates suggest the parallel corridors add 50--100\% to the addressable
market. Total coast-to-coast air-to-truck substitution savings, including all major
corridors, may reach \$12--16 billion per year at full adoption --- comparable in
magnitude to the annual benefit stream estimated for the entire I2.0 managed lane
program \citep{ROUTE_E1}.

\subsection{A Note on Second-Order Effects}

The \$8.2 billion in direct freight cost savings understates the total economic benefit
for three reasons.

\textit{Lower final goods prices.} Freight costs are embedded in the prices consumers pay.
A 90\% reduction in shipping cost for California strawberries does not translate to a
90\% reduction in the retail price of strawberries --- but it does translate to something.
Lower logistics costs in competitive markets are partly passed through to consumers.

\textit{Expanded producer markets.} The California strawberry farmer who currently cannot
reach New York fresh gains a new market. This is not a freight cost saving; it is new
revenue for the producer. We quantify this separately in Section~\ref{sec:fresh}.

\textit{Inventory carrying cost reduction.} Products that must air-freight because truck
transit is too slow also tend to be held at low inventory levels (because of the air
freight cost). With 48-hour truck, safety stock can be higher and carrying cost per unit
can be lower simultaneously. This inventory optimization benefit is real but difficult to
estimate without shipper-level data; we note it as an unmeasured upside.

The combined effect of these three second-order mechanisms --- consumer price pass-through,
producer market access, and inventory optimization --- suggests that the \$8.2 billion
headline savings figure understates the full economic value of air-to-truck substitution
by a meaningful margin. Sections~\ref{sec:fresh} through~\ref{sec:ecommerce} develop
the producer market access and inventory cases in detail for the produce, e-commerce,
and pharmaceutical sectors. In each, the direct freight cost saving is the floor of the
economic benefit, not the ceiling.
